\section{Hipótesis}\label{cap:hipotesis}

La hipótesis expresa una respuesta tentativa al problema formulado, que será contrastada con los datos de la línea base retrospectiva del salón y de las semanas de prueba piloto.


\textbf{Hipótesis de investigación (Hi):} Si se implementa el sistema web móvil Lush Beauty para la gestión de caja y nómina, entonces disminuyen el tiempo de liquidación semanal y la tasa de errores de cálculo de comisiones en el salón Lush Beauty de la ciudad de Santa Cruz de la Sierra durante el piloto de enero de 2027.


\textbf{Hipótesis nula ($H_0$):} La implementación del sistema web móvil Lush Beauty no tiene un efecto significativo en el tiempo de liquidación semanal ni en la tasa de errores de cálculo de comisiones del salón.


\textbf{Hipótesis alternativa (Ha):} La implementación del sistema web móvil Lush Beauty tiene un efecto significativo en el tiempo de liquidación semanal y/o en la tasa de errores de cálculo de comisiones.


\subsection{Hipótesis para el soporte a decisiones}

La ampliación hacia el soporte para decisiones incorpora, además, las siguientes hipótesis de trabajo, referidas a tareas estructuradas de comparación financiera en el entorno piloto.

Hipótesis de investigación (Hi): El uso del prototipo Lush Beauty Smart DSS permitirá a las participantes completar tareas estructuradas de comparación financiera con menos tiempo y mayor trazabilidad de criterios que con el procedimiento manual basado en hojas de cálculo, en el entorno piloto definido para el estudio.


Hipótesis nula ($H_0$): En las tareas estructuradas del piloto, no se observará una mejora en el tiempo de comparación de alternativas ni en la trazabilidad de criterios respecto del procedimiento manual.


Hipótesis alternativa operativa (Ha): Se observará mejora en al menos uno de los indicadores de tiempo o trazabilidad; cualquier diferencia deberá interpretarse con prudencia por el tamaño reducido de la muestra y la ausencia de grupo de control.

